\labsection{3}{P2P loan-book economics and diversification}{sec:lab03-p2p}

\begin{labproject}{Price credit risk, then prove diversification does what the chapter claims}
\textbf{Coverage.} Chapter~\ref{ch:p2p-crowdfunding}.\\
\textbf{Goal.} Compute the return an investor actually receives on a P2P loan book, and demonstrate numerically why concentration is the risk that ruins retail lenders.

\courselabflow{Inspect loans}{Estimate loss}{Compare grades}{Simulate portfolios}



\end{labproject}

\textbf{Complete worked example.} This runnable fixture uses synthetic inputs. First reproduce it, then adapt the data or ratings with documented evidence. All code below belongs to one file.
\begin{lstlisting}[style=utpterminal]
python all_experiments/lab03_p2p_loan_book.py
\end{lstlisting}

\noindent\textbf{Setup and output folder.}
\lstinputlisting[style=utppython,firstline=1,lastline=10]{all_experiments/lab03_p2p_loan_book.py}

\noindent\textbf{Create loans with transparent grade assumptions.}
\lstinputlisting[style=utppython,firstline=11,lastline=24]{all_experiments/lab03_p2p_loan_book.py}

\noindent\textbf{Use one consistent cash-flow model.}
\courselabtrace{From coupon to a classroom net-yield estimate}{Credit loss}{Default probability times loss given default}{Fees}{Deduct the explicitly stated platform-fee assumption}{Diversification}{Compare dispersion; correlation can limit the benefit}



\textbf{Use the chapter while coding.} Revisit \hyperref[ch:p2p-crowdfunding]{the earlier chapter explanation}; match its operation to the variables and state changes below.

\textbf{Follow the supplied implementation.} Return and loss are different columns with different meanings. Follow np.where to see which return branch applies, then inspect one actual default row. The portfolio bootstrap samples these realised loan returns independently; it does not model correlated default shocks.
\courselabtrace{Trace the implemented loan return convention}{1. Input / state}{Each synthetic loan has interest\_rate, recovery\_rate and defaulted; equal principal is assumed and fee = .01.}{2. Operation}{Non-default return is interest\_rate minus fee. Default return is recovery\_rate minus 1 minus fee.}{3. Check}{loss is defaulted times (1 - recovery\_rate); group means summarise by grade.}


\lstinputlisting[style=utppython,firstline=25,lastline=36]{all_experiments/lab03_p2p_loan_book.py}

\noindent\textbf{Resample equal-weight portfolios.}
\lstinputlisting[style=utppython,firstline=37,lastline=48]{all_experiments/lab03_p2p_loan_book.py}
\begin{labvisual}{Synthetic IID portfolios: dispersion of net returns}
\begin{labgraphic}\begin{tikzpicture}[font=\small]
\node[anchor=east,font=\scriptsize] at (-.2,0.0) {1 loans};
\fill[UTPPaleBlue] (0,-0.16) rectangle (7,0.16);
\fill[UTPTeal] (0,-0.16) rectangle (7.0000,0.16);
\node[anchor=west,font=\scriptsize] at (7.15,0.0) {28.0\% SD};
\node[anchor=east,font=\scriptsize] at (-.2,-0.55) {5 loans};
\fill[UTPPaleBlue] (0,-0.7100000000000001) rectangle (7,-0.39);
\fill[UTPTeal] (0,-0.7100000000000001) rectangle (3.2201,-0.39);
\node[anchor=west,font=\scriptsize] at (7.15,-0.55) {12.9\% SD};
\node[anchor=east,font=\scriptsize] at (-.2,-1.1) {20 loans};
\fill[UTPPaleBlue] (0,-1.26) rectangle (7,-0.9400000000000001);
\fill[UTPTeal] (0,-1.26) rectangle (1.5897,-0.9400000000000001);
\node[anchor=west,font=\scriptsize] at (7.15,-1.1) {6.4\% SD};
\node[anchor=east,font=\scriptsize] at (-.2,-1.6500000000000001) {100 loans};
\fill[UTPPaleBlue] (0,-1.81) rectangle (7,-1.4900000000000002);
\fill[UTPTeal] (0,-1.81) rectangle (0.7113,-1.4900000000000002);
\node[anchor=west,font=\scriptsize] at (7.15,-1.6500000000000001) {2.8\% SD};
\node[anchor=east,font=\scriptsize] at (-.2,-2.2) {500 loans};
\fill[UTPPaleBlue] (0,-2.3600000000000003) rectangle (7,-2.04);
\fill[UTPTeal] (0,-2.3600000000000003) rectangle (0.3128,-2.04);
\node[anchor=west,font=\scriptsize] at (7.15,-2.2) {1.2\% SD};
\end{tikzpicture}\end{labgraphic}\end{labvisual}

\textbf{Check the run.} The following output was captured from the supplied script; it describes the synthetic fixture only.
\lstinputlisting[style=utpterminal]{labs/outputs/lab03.txt}
% Reused from chapters/03-p2p-and-crowdfunding.tex
\begin{labvisual}{Recall the chapter model}
\begin{labgraphic}
\begin{tikzpicture}[
  p/.style={draw=UTPBlue,fill=UTPPaleBlue,rounded corners=2mm,text width=2.6cm,
            minimum height=1.15cm,align=center,font=\small\bfseries\color{UTPNavy}},
  mid/.style={draw=UTPNavy,fill=white,rounded corners=2mm,text width=3.4cm,
              minimum height=1.4cm,align=center,font=\small\bfseries\color{UTPNavy}},
  arr/.style={-{Stealth[length=2.1mm]},thick,draw=UTPTeal},
  head/.style={font=\small\bfseries\color{UTPNavy}},
  risk/.style={font=\scriptsize\color{UTPRed},align=center,text width=4.4cm},
  ring/.style={draw=UTPRed,dashed,rounded corners=2.5mm,line width=0.8pt}]

% ---- traditional intermediation ---------------------------------
\node[head] at (0,1.55) {Traditional intermediation};
\node[p]   (s1) at (-4.6,0) {Savers};
\node[mid] (k1) at (0,0)    {Bank\\own balance sheet};
\node[p]   (b1) at (4.6,0)  {Borrowers};
\draw[arr] (s1)--node[above,font=\scriptsize]{deposits}(k1);
\draw[arr] (k1)--node[above,font=\scriptsize]{loans}(b1);
\draw[ring] (-1.85,-0.78) rectangle (1.85,0.78);
\node[risk] at (0,-1.20) {a default is absorbed here, behind reserves and deposit insurance};

% ---- marketplace funding ----------------------------------------
\node[head] at (0,-2.55) {Marketplace funding};
\node[p]   (s2) at (-4.6,-4.05) {Lenders};
\node[mid] (k2) at (0,-4.05)    {Platform\\matches, scores, services};
\node[p]   (b2) at (4.6,-4.05)  {Borrowers};
\draw[arr] (s2)--node[above,font=\scriptsize]{capital}(k2);
\draw[arr] (k2)--node[above,font=\scriptsize]{loans}(b2);
\draw[ring] (-6.0,-4.83) rectangle (-3.2,-3.27);
\node[risk] at (-4.6,-5.25) {a default is absorbed here, by the individual lender};

\end{tikzpicture}
\end{labgraphic}\end{labvisual}


\begin{verification}
\textbf{Explain and extend your result.} Compare portfolio sizes using mean return and downside quantiles. Explain why a narrower simulated distribution is conditional on the independence assumption.

\begin{itemize}\item The script runs from a fresh process and creates the named CSV or report files.\item Every rate uses a stated population and denominator; every model assumption is identified.\item Synthetic outputs are labelled and are not presented as observations of a real institution.\item Change one input, rerun, and explain the direction of the result.\end{itemize}\end{verification}
\begin{labdeliverable}Submit the complete script, generated evidence, and a one-page interpretation linking the results to the chapter model. State the assumptions and one limitation.\end{labdeliverable}
